\documentclass[12pt,a4paper]{article}
\usepackage[utf8]{inputenc}
\usepackage[T1]{fontenc}
\usepackage[margin=2.5cm]{geometry}
\usepackage{graphicx}
\usepackage{listings}
\usepackage{xcolor}
\usepackage{caption}
\usepackage{booktabs}
\usepackage{array}
\newcommand{\studentgroup}{FAF-242}
\definecolor{termbg}{rgb}{0.96,0.96,0.96}
\lstset{backgroundcolor=\color{termbg},basicstyle=\ttfamily\footnotesize,breaklines=true,frame=single,showstringspaces=false,columns=fullflexible,keepspaces=true}
\begin{document}
\begin{titlepage}
    \centering
    \includegraphics[width=0.25\textwidth]{utm_logo.png}\\[0.6cm]
    {\large \textbf{UNIVERSITATEA TEHNICĂ A MOLDOVEI}}\\[0.2cm]
    {\large Facultatea Calculatoare, Informatică și Microelectronică}\\[0.2cm]
    {\large Departamentul Ingineria Software și Automatică}\\[2.5cm]
    {\Large \textbf{LUCRAREA DE LABORATOR NR. 2}}\\[0.4cm]
    {\LARGE \textbf{Utilitarul \texttt{sleep} în xv6}}\\[0.4cm]
    {\large Disciplina: Sisteme de Operare}\\[2.5cm]
    \begin{minipage}{0.45\textwidth}\flushleft\textbf{Done by:}\\st. gr. \studentgroup\\Daniel-Ioan \textsc{Strunga}\end{minipage}\hfill
    \begin{minipage}{0.50\textwidth}\flushright\textbf{Checked by:}\\asist. univ. Patricia \textsc{Reițman}\\asist. univ. Marius \textsc{Dumitrașcu}\end{minipage}
    \vfill{\large Chișinău --- 2026}
\end{titlepage}
\newpage

\section{Scopul lucrării}
Compilarea și rularea sistemului xv6-riscv în QEMU și adăugarea unui program utilizator nou, \texttt{sleep}, care primește numărul de tick-uri ca argument și apelează apelul de sistem corespunzător din kernel.

\section{Partea 1 și 2 --- rulare și explorare xv6}
\subsection{Comenzi rulate în xv6 (QEMU)}
\begin{lstlisting}[language=bash]
$ make qemu
$ ls
$ cat README
$ echo hello xv6
hello xv6
$ ls | grep c
cat            2 3 36288
echo           2 4 35152
wc             2 19 37240
sync           2 24 34600
console        3 25 0
$ wc README
48 336 2441 README
\end{lstlisting}

\subsection{Răspunsuri (Partea 2)}
\begin{enumerate}
    \item \textbf{Trei programe incluse în xv6:} \texttt{ls}, \texttt{cat}, \texttt{echo} (plus \texttt{grep}, \texttt{wc}, \texttt{sh}).
    \item \textbf{Două funcționalități ale SO necesare pentru pipe-uri:}
    \begin{itemize}
        \item procese multiple care rulează concurent: shell-ul creează câte un proces pentru fiecare parte a pipe-ului (\texttt{fork} + \texttt{exec}), iar planificatorul le alternează;
        \item comunicarea între procese prin descriptori de fișiere: kernelul creează pipe-ul (\texttt{pipe()}), iar shell-ul mută capetele lui pe descriptorii 0 și 1 (\texttt{close} + \texttt{dup}).
    \end{itemize}
    \item \textbf{Shell-ul xv6 față de shell-ul Linux:} shell-ul xv6 este minimal: execută programe, are redirectări (\texttt{<}, \texttt{>}), pipe-uri, \texttt{;} și \texttt{\&}. Îi lipsesc istoricul, completarea cu Tab, variabilele de mediu, globbing-ul și \texttt{PATH}, pe care Bash le are.
\end{enumerate}

\section{Partea 3 --- analiza codului sursă}
\subsection{Comenzi rulate în \texttt{\textasciitilde/xv6-riscv}}
\begin{lstlisting}[language=bash]
$ sed -n "1,30p" user/cat.c
$ cat user/user.h
$ grep -n "sys_read\|sys_write" kernel/sysfile.c | head
69:sys_read(void)
83:sys_write(void)
$ grep -n -o "\b\(read\|write\|open\|close\|exit\)(" user/cat.c | sort -t: -k2 -u
40:close(
15:exit(
35:open(
12:read(
13:write(
\end{lstlisting}

\subsection{Răspunsuri (Partea 3)}
\begin{enumerate}
    \item \textbf{Apeluri de sistem folosite în \texttt{user/cat.c}:}
    \begin{itemize}
        \item \texttt{read()} --- cere kernelului octeți dintr-un descriptor de fișier, puși într-un buffer al procesului;
        \item \texttt{write()} --- cere kernelului să scrie octeții din buffer într-un descriptor (1 = stdout);
        \item \texttt{open()} / \texttt{close()} --- deschid și eliberează fișierul primit ca argument;
        \item \texttt{exit()} --- termină procesul și eliberează resursele lui.
    \end{itemize}
    \item \textbf{Unde este implementat \texttt{sys\_read}:} în \texttt{kernel/sysfile.c}, la linia 69 (\texttt{sys\_write} la linia 83).
    \item \textbf{Diferența dintre \texttt{kernel/} și \texttt{user/}:} codul din \texttt{kernel/} rulează în modul privilegiat (supervisor mode), cu acces direct la hardware și la memoria fizică. Codul din \texttt{user/} rulează în user mode, în spațiul de adrese al procesului, și ajunge la resurse doar prin apeluri de sistem (\texttt{ecall}).
\end{enumerate}

\section{Partea 4 --- programul \texttt{sleep}}
\subsection{Codul programului (\texttt{user/sleep.c})}
{\sloppy Pe versiunea curentă de xv6 apelul de sistem care suspendă procesul se numește \texttt{pause(int)}; numele \texttt{sleep} a rămas doar în mecanismul intern de adormire a firelor din kernel (\texttt{sleep\_prepare()} în \texttt{kernel/proc.c}), deci un program care apelează \texttt{sleep(n)} nu se compilează.\par}
\begin{lstlisting}[language=C]
#include "kernel/types.h"
#include "user/user.h"

int
main(int argc, char *argv[])
{
  if(argc != 2){
    fprintf(2, "usage: sleep <ticks>\n");
    exit(1);
  }
  pause(atoi(argv[1])); // was sleep() before Aug 2025
  exit(0);
}
\end{lstlisting}

\subsection{Integrarea în \texttt{Makefile}}
În lista \texttt{UPROGS} s-a adăugat linia \verb|$U/_sleep\|. Regula generală \verb|$U/_%: $U/%.o ...| compilează \texttt{sleep.c}, îl leagă cu \texttt{ulib.o}, \texttt{usys.o}, \texttt{printf.o}, \texttt{umalloc.o}, iar \texttt{mkfs} copiază binarul \texttt{\_sleep} în imaginea \texttt{fs.img} sub numele \texttt{sleep}.

\subsection{Rulare în QEMU}
\begin{figure}[h!]
\begin{lstlisting}
$ sleep
usage: sleep <ticks>
$ echo before; sleep 100; echo after
before
after
$ ls
...
sleep          2 5 34904
...
\end{lstlisting}
\caption{Figure 2.1 -- \texttt{sleep} fără argument, cu 100 de tick-uri și prezența în sistemul de fișiere}
\end{figure}

\subsection{Răspunsuri (Observe)}
\begin{enumerate}
    \item \textbf{Ce unitate de timp are argumentul?} Tick-uri ale timer-ului kernelului (în xv6 pe QEMU un tick este aproximativ 0,1 s), nu secunde. Între \texttt{before} și \texttt{after} trece vizibil timp, pentru că \texttt{sleep 100} blochează procesul circa 10 s.
    \item \textbf{Ce se întâmplă fără argument?} Programul verifică \texttt{argc != 2}, scrie mesajul pe descriptorul 2 (stderr) și iese cu cod 1, ca în figura 2.1.
    \item \textbf{De ce \texttt{sleep} este program separat și nu comandă internă?} Nu modifică starea shell-ului (spre deosebire de \texttt{cd}), așa că poate rula într-un proces copil după \texttt{fork}/\texttt{exec}.
    \item \textbf{Cum ajunge apelul în kernel?} \texttt{pause()} este un stub generat de \texttt{usys.pl} care pune numărul apelului în \texttt{a7} și execută \texttt{ecall}; kernelul îl primește în \texttt{syscall()} și apelează \texttt{sys\_pause}, care adoarme procesul pe canalul \texttt{ticks} până când trec tick-urile cerute.
    \item \textbf{Ce face \texttt{atoi}?} Convertește șirul din \texttt{argv[1]} în întreg; o valoare nenumerică dă 0 și programul iese imediat.
\end{enumerate}

\section{Concluzii}
Am compilat xv6, am adăugat programul \texttt{sleep} folosind \texttt{pause()} (apelul de sistem actual), l-am înregistrat în \texttt{UPROGS} și am verificat în QEMU cazul de eroare și cazul normal. Lucrarea a clarificat drumul complet: program utilizator, stub \texttt{usys}, \texttt{ecall}, tabelul de apeluri de sistem din kernel.
\end{document}
